% Source-checked bibliography, audit date 15 November 2026.
% S checks establish the stated source scope; they do not substitute for proof checking.
\begin{thebibliography}{FKN-e}
\bibitem[P]{P} OpenAI, \emph{A polynomial-time algorithm for three-machine unit-job scheduling},
OpenAI Math Release preprint, 24 September 2026,
\url{https://github.com/openai/math/tree/main/preprints/A-polynomial-time-algorithm-for-three-machine-unit-job-scheduling-September-24-2026}.
Opened at repository commit \texttt{4eaee87548250b2682affda7d05d947c05554bae};
the thirteen attached source files are byte-identical to that version. [S-P]
\bibitem[F]{F} OpenAI Math Release, formalization scope file \texttt{lean/docs/124.md},
\url{https://github.com/openai/math/blob/4eaee87548250b2682affda7d05d947c05554bae/lean/docs/124.md}.
Scope, challenge declaration and actual \texttt{Scheduling/Main.lean} proof body opened;
the comparator's \texttt{sorry} is explicitly distinguished from the separate proof source.
No Lean build replayed in this audit. [S-F]
\bibitem[D]{D} drbh, \emph{scheduling-proof-exploration},
\url{https://github.com/drbh/scheduling-proof-exploration/tree/bbd30137258b15d8c5caedd035d72e8e8a875dbd},
commit \texttt{bbd30137258b15d8c5caedd035d72e8e8a875dbd}, 8 October 2026.
The source declares and proves the full four-machine release theorem, with
explicit job encoding, exact optimization and deadline output, and a fixed
multitape polynomial bound; the reported Lean build was not replayed here. [S-D]
\bibitem[N]{N} J.~Nederlof, C.~M.~F. Swennenhuis and K.~W\k{e}grzycki,
A subexponential time algorithm for makespan scheduling of unit jobs with precedence constraints,
\emph{Proceedings of the 2025 Annual ACM--SIAM Symposium on Discrete Algorithms (SODA)},
535--552 (2025),
\href{https://doi.org/10.1137/1.9781611978322.16}{doi:10.1137/1.9781611978322.16};
\url{https://arxiv.org/abs/2312.03495}, version 1, 6 December 2023. [S-N]
\bibitem[A]{A} S.~Das and A.~Wiese,
A simpler QPTAS for scheduling jobs with precedence constraints,
\emph{30th Annual European Symposium on Algorithms (ESA 2022)},
LIPIcs 244, 40:1--40:11 (2022),
\href{https://doi.org/10.4230/LIPIcs.ESA.2022.40}{doi:10.4230/LIPIcs.ESA.2022.40};
\url{https://arxiv.org/abs/2501.09091}, version 1, 15 January 2025. [S-A]
\bibitem[FKN]{FKN} M.~Fujii, T.~Kasami and K.~Ninomiya,
Optimal sequencing of two equivalent processors,
\emph{SIAM Journal on Applied Mathematics} 17(4), 784--789 (1969),
\href{https://doi.org/10.1137/0117070}{doi:10.1137/0117070}.
Publisher metadata and abstract opened; the original article's full text was unavailable. [S-FKN]
\bibitem[FKN-e]{FKNe} M.~Fujii, T.~Kasami and K.~Ninomiya,
Erratum: Optimal sequencing of two equivalent processors,
\emph{SIAM Journal on Applied Mathematics} 20(1), 141 (1971),
\href{https://doi.org/10.1137/0120018}{doi:10.1137/0120018}.
Erratum full text opened. [S-FKN-e]
\bibitem[CG]{CG} E.~G. Coffman, Jr. and R.~L. Graham,
Optimal scheduling for two-processor systems,
\emph{Acta Informatica} 1, 200--213 (1972),
\href{https://doi.org/10.1007/BF00288685}{doi:10.1007/BF00288685}. [S-CG]
\bibitem[U]{U} J.~D. Ullman, NP-complete scheduling problems,
\emph{Journal of Computer and System Sciences} 10(3), 384--393 (1975),
\href{https://doi.org/10.1016/S0022-0000(75)80008-0}{doi:10.1016/S0022-0000(75)80008-0}. [S-U]
\bibitem[GJ]{GJ} M.~R. Garey and D.~S. Johnson,
\emph{Computers and Intractability: A Guide to the Theory of NP-Completeness},
W.~H. Freeman, San Francisco (1979), Appendix A13, OPEN8, p.~287;
also Appendix A5.2, SS9, p.~239. [S-GJ]
\bibitem[L]{L} Open Mathematics Courses,
\emph{Scheduling unit jobs on three machines}, lessons 1--11,
\url{https://kokunoyumeto.github.io/open-math-courses-public/courses/three-machine-unit-job-scheduling/}.
Course byline: Claude Opus 5.5, effort not recorded. Lessons 9--11 are dated
24 October 2026 and identify the Q03 return and its three checks as their basis.
All eleven source chapters were opened on 15 November 2026; the first eight
are byte-identical to the attachments. [S-L]
\bibitem[J]{J} K.~Jansen, \emph{Precedence constraint scheduling on a fixed number of machines},
\url{https://arxiv.org/abs/2611.09392}, version 1, 6 November 2026.
Theorem 1.1 and Lemma 7.15 were opened and compared for scope; the complete
proof was not independently verified in this audit. [S-J]
\bibitem[G]{G} D.~N. Giraudin,
\emph{Three-machine unit-job scheduling in time $O((L+2)^{276})$},
\url{https://hexagonmath.org/pdf/2611.00060v1}, version 1, November 2026.
Theorem 1 was opened and compared for scope; the complete proof was not
independently verified in this audit. [S-G]
\end{thebibliography}
